% ============================================================
% Section 1.2: Problem Statement
% ============================================================
\section{Problem Statement}
\label{sec:problem_statement}

Individuals with severe motor paralysis face extreme social isolation and total reliance on caregivers due to the lack of functional communication and mobility tools. Although various assistive technologies currently exist, they present major technical and practical limitations. First, invasive Brain-Computer Interfaces (BCIs) provide direct neural control but carry significant surgical risks, tissue infection hazards, and high financial costs. Second, Video-Oculography (VOG) camera-based eye-tracking alternatives are highly sensitive to ambient lighting variations, require continuous spatial calibration, and demand high computational processing power. Third, system fragmentation remains a major barrier, as most commercially available devices treat communication and physical locomotion as separate challenges, requiring independent and expensive standalone platforms for each function.


Consequently, an engineering need exists for a non-invasive, low-cost HMI with reliable threshold detection. Such a system must utilize simplified biopotential acquisition hardware to deliver real-time control over both digital text communication and motorized mobility platforms within a single interface.

This project presents the design and implementation of an assistive technology system that integrates a virtual keyboard for text-based communication and a prototype motorized wheelchair controlled via real-time EOG thresholding. This framework addresses mobility and communication requirements within a single, unified platform.

